\begin{statementbox}
\textbf{\textcolor{CStatement}{条件与结论}}
设 $f(x)=a_nx^n+\cdots+a_1x+a_0$，其中 $n\ge1$，各系数为整数，$a_n\ne0$、$a_0\ne0$。若最简分数 $p/q$ 是根，且 $q>0$、$p,q$ 互质，则
\[
  \boxed{p\mid a_0,\qquad q\mid a_n.}
\]
这里 $p\mid a_0$ 表示 $a_0$ 是 $p$ 的整数倍。\textbf{分子看常数项，分母看最高次项系数}；这只是必要条件，候选值仍须代入检验。

\begin{center}
\begin{tikzpicture}[x=1cm,y=1cm,font=\small]
\node[font=\small\bfseries] at (5.5,3.55) {$f(x)=2x^2+x-1$：四个候选，两个真根};
\node at (1.25,2.65) {$q\backslash p$};
\node at (4.15,2.65) {$-1\mid(-1)$};
\node at (8.05,2.65) {$1\mid(-1)$};
\fill[CExample!10] (2.3,1.2) rectangle (6.2,2.2);
\fill[CExample!10] (6.2,0.15) rectangle (10.1,1.2);
\draw[CStatement!45] (2.3,0.15) rectangle (10.1,2.2);
\draw[CStatement!45] (6.2,0.15)--(6.2,2.2);
\draw[CStatement!45] (2.3,1.2)--(10.1,1.2);
\node at (1.25,1.7) {$1\mid2$};
\node at (1.25,0.65) {$2\mid2$};
\node[text=CExample] at (4.15,1.7) {$x=-1,\ f(x)=0$};
\node at (8.05,1.7) {$x=1,\ f(x)=2$};
\node at (4.15,0.65) {$x=-\frac12,\ f(x)=-1$};
\node[text=CExample] at (8.05,0.65) {$x=\frac12,\ f(x)=0$};
\node at (5.5,-0.45) {行选分母、列选分子；绿色格才对应验出的根。};
\end{tikzpicture}
\end{center}


\textbf{\textcolor{CStatement}{验根后才能提因式}}
若 $f(p/q)=0$，则
\[
  \boxed{f(x)=(qx-p)g(x),\qquad g(x)\text{ 为整数系数多项式}.}
\]
商式次数为 $n-1$；继续分解或求根，并检查重复因式。

\medskip
\textbf{\textcolor{CStatement}{两个先行判断}}
\begin{itemize}[leftmargin=*,itemsep=0.4em]
\item 若 $a_n=\pm1$，则 $q=1$，有理根只能是常数项的整数约数。
\item 若 $a_0=0$，先提取所有公因式 $x$，记下零根及其重数，再对剩余的非常数多项式应用定理。
\end{itemize}
\end{statementbox}
\clearpage
